\documentclass[]{article}

\usepackage{amsmath}
\usepackage{multirow}
\usepackage{natbib}
\usepackage{graphicx} 


\begin{document}

% Auto-generated figures file
% \usepackage{graphicx}

\begin{figure}[h!]
    \centering
    \includegraphics[width=0.5\textwidth]{../input_files/plots/step_3_satellite_radial_profiles_1_20260414_163828.png}
    \caption{Radial density profiles of satellite galaxies, shown for four distinct host halo mass bins. The empirical profiles from the synthetic catalog (black points with error bars) are compared against the true data-generating model (blue dashed line) and the best-fit exponential model recovered via maximum likelihood estimation (red solid line). The recovered model provides a strong fit to the data across all mass scales, but exhibits a systematic offset above the true model, visually representing the slight positive bias in the recovered Halo Occupation Distribution parameters, $M\_{sat}$ and $\alpha\_{sat}$.}
    \label{fig:image0}
\end{figure}

\begin{figure}[h!]
    \centering
    \includegraphics[width=0.5\textwidth]{../input_files/plots/step_4_marked_correlation_1776184816.png}
    \caption{The luminosity-marked correlation function ratio, $M(r)/M\_{shuffled}(r)$, for the 1-halo ($r < 5$ Mpc/h, left panel) and 2-halo ($r > 10$ Mpc/h, right panel) regimes. In the 1-halo regime, the ratio is significantly greater than unity at small separations, which demonstrates strong spatial-luminosity segregation arising from the pairing of systematically brighter central galaxies with fainter satellites. In the 2-halo regime, the ratio is consistent with unity, confirming that galaxy luminosities are spatially uncorrelated on large scales where galaxy pairs reside in distinct, independent halos.}
    \label{fig:image1}
\end{figure}

\begin{figure}[h!]
    \centering
    \includegraphics[width=0.5\textwidth]{../input_files/plots/step_5_model_comparison_residuals_1_20260414_164420.png}
    \caption{Left: The recovered concentration parameter of the satellite radial profile, $\alpha(M)$, as a function of host halo mass, $M\_{vir}$. The analysis shows a decisive statistical preference for the mass-dependent model (red solid line), which indicates that more massive halos host more concentrated satellite distributions, over the constant-concentration model (blue dashed line). Right: The fractional residual, $(O-E)/E$, between the observed (O) and expected (E) satellite counts for the preferred mass-dependent model in the transition region between the 1-halo and 2-halo regimes. The dramatic increase in residuals at distances beyond $\sim$7 Mpc/h demonstrates the fundamental breakdown of the isolated 1-halo model, which severely underpredicts galaxy counts by neglecting the 2-halo contribution from clustered neighboring halos.}
    \label{fig:image2}
\end{figure}



\end{document}
                